% ingredients.tex
\usepackage{xfp} % Modern high-precision floating point calculation engine
\usepackage{tcolorbox}
\usepackage{booktabs}

% --- LATEX3 KEY-VALUE SETUP & MATH ENGINE ---
\ExplSyntaxOn

% Define the nutrition key storage fields
\keys_define:nn { ing }
  {
    name    .tl_set:N  = \l_ing_name_tl,
    kcal    .fp_set:N  = \l_ing_kcal_fp,
    protein .fp_set:N  = \l_ing_protein_fp,
    carbs   .fp_set:N  = \l_ing_carbs_fp,
    fat     .fp_set:N  = \l_ing_fat_fp,
    satfat  .fp_set:N  = \l_ing_satfat_fp,
    sugar   .fp_set:N  = \l_ing_sugar_fp,
    fiber   .fp_set:N  = \l_ing_fiber_fp,
    salt    .fp_set:N  = \l_ing_salt_fp,
  }

% Native LaTeX3 definition of the print macro to prevent parameter bugs
\cs_new:Npn \printValue #1 
  {
    \fpeval{round(#1, 2)}
  }

% Native LaTeX3 initialization of calculation registers
\cs_new:Npn \initRecipeCalculator 
  {
    \gdef\totalKcal{0}
    \gdef\totalProtein{0}
    \gdef\totalCarbs{0}
    \gdef\totalFat{0}
    \gdef\totalSatFat{0}
    \gdef\totalSugar{0}
    \gdef\totalFiber{0}
    \gdef\totalSalt{0}
  }

% Native LaTeX3 definition to save ingredients into the global data cache
\cs_new:Npn \defIngredient #1 #2
  {
    \expandafter\gdef\csname ingdata@#1\endcsname{#2}
  }

% Core logic to add an ingredient, evaluate math and print layout columns
\newcommand{\addIngredient}[2]{
    % Reset variables to default 0 values safely
    \keys_set:nn { ing } { name={#1}, kcal=0, protein=0, carbs=0, fat=0, satfat=0, sugar=0, fiber=0, salt=0 }
    % Load and parse the stored key-value list for this specific ID
    \keys_set:nv { ing } { ingdata@#1 }
    %
    % Export the token list name safely to a standard global macro so tabular can read it
    \tl_gset_eq:NN \g_printed_name_tl \l_ing_name_tl
    \expandafter\gdef\expandafter\printedName\expandafter{\g_printed_name_tl}
    %
    % Accumulate values using floating point evaluation
    \xdef\totalKcal{\fpeval{\totalKcal + (#2 * \l_ing_kcal_fp / 100)}}
    \xdef\totalProtein{\fpeval{\totalProtein + (#2 * \l_ing_protein_fp / 100)}}
    \xdef\totalCarbs{\fpeval{\totalCarbs + (#2 * \l_ing_carbs_fp / 100)}}
    \xdef\totalFat{\fpeval{\totalFat + (#2 * \l_ing_fat_fp / 100)}}
    \xdef\totalSatFat{\fpeval{\totalSatFat + (#2 * \l_ing_satfat_fp / 100)}}
    \xdef\totalSugar{\fpeval{\totalSugar + (#2 * \l_ing_sugar_fp / 100)}}
    \xdef\totalFiber{\fpeval{\totalFiber + (#2 * \l_ing_fiber_fp / 100)}}
    \xdef\totalSalt{\fpeval{\totalSalt + (#2 * \l_ing_salt_fp / 100)}}
    %
    % Output table row with exactly 2 decimal places
    \fpeval{round(#2, 2)} & g & \printedName \\
}
\ExplSyntaxOff

% --- NUTRITION FACTS PANEL GRAPHIC ---
\newcommand{\printNutritionFacts}{%
    \begin{tcolorbox}[colback=white, colframe=black, sharp corners, boxrule=1pt, width=6.0cm, title=\textbf{Nutrition Facts (total)}, coltitle=black, colbacktitle=white, attach title to upper=\vspace{0.4em}]
        \small
        \hrule width \linewidth height 2pt \vspace{0.3em}
        \begin{tabular}{@{} l r l @{}}
        \noindent Energy        & \printValue{\totalKcal}    & kcal \\
        \noindent Fat           & \printValue{\totalFat}     & g    \\
        \noindent Saturated Fat & \printValue{\totalSatFat}  & g    \\
        \noindent Carbohydrate  & \printValue{\totalCarbs}   & g    \\
        \noindent Fiber         & \printValue{\totalFiber}   & g    \\
        \noindent Sugars        & \printValue{\totalSugar}   & g    \\
        \noindent Protein       & \printValue{\totalProtein} & g    \\
        \noindent Sodium / Salt & \printValue{\totalSalt}    & g    \\
        \end{tabular}
    \end{tcolorbox}%
}

% Updated box style that dynamically adapts its width to its surrounding container
\newtcolorbox{ingredientsboxfull}[1][\linewidth]{
    sharp corners, colback=gray!3, colframe=gray!30, boxrule=0.5pt,
    width=#1, top=10pt, bottom=10pt, left=10pt, right=10pt
}

% =======================================================
% CENTRAL DATA STORAGE: KEY-VALUE FORMAT (Values per 100g)
% =======================================================
\defIngredient{ryeflour}{name={Rye Whole Grain Flour}, kcal=339, protein=9, carbs=62.4, fat=2.3, satfat=0.3, sugar=5.64, fiber=14.8, salt=0.00}
\defIngredient{water}{name={Water}, kcal=0, protein=0, carbs=0, fat=0, satfat=0, sugar=0, fiber=0, salt=0}
\defIngredient{waterBoiling}{name={Water (boiling)}, kcal=0, protein=0, carbs=0, fat=0, satfat=0, sugar=0, fiber=0, salt=0}
\defIngredient{salt}{name={Sea Salt}, kcal=0, protein=0, carbs=0, fat=0, satfat=0, sugar=0, fiber=0, salt=99.0}
\defIngredient{sourdough}{name={Sourdough Starter}, kcal=169.5, protein=4.5, carbs=31.2, fat=1.15, satfat=0.15, sugar=2.82, fiber=7.4, salt=0.00}
\defIngredient{coarseRyeMeal75}{name={Coarse rye meal (grind size 7.5)}, kcal=339, protein=9, carbs=62.4, fat=2.3, satfat=0.3, sugar=5.64, fiber=14.8, salt=0.00}
\defIngredient{coarseRyeMeal75FootNote}{name={Coarse rye meal (grind size 7.5)\footnote{Setting on "hawos M001 M1 Grinder, 1.3 Kg, 360 W, Wood"}}, kcal=339, protein=9, carbs=62.4, fat=2.3, satfat=0.3, sugar=5.64, fiber=14.8, salt=0.00}
\defIngredient{coarseRyeMeal3}{name={Coarse rye meal (grind size 3)}, kcal=339, protein=9, carbs=62.4, fat=2.3, satfat=0.3, sugar=5.64, fiber=14.8, salt=0.00}

